%!TEX program=xelatex
%===============================================================================
%  实验报告 LaTeX 模板
%  基于 SDU、SimpleLabReport、NUSRI、NUIST 等优秀模板的学习与整合
%  适用于算法与数据分析类实验课程
%===============================================================================
\documentclass[12pt,a4paper,UTF8]{article}

%-----------------------------------------------------------------------------
%  宏包加载
%-----------------------------------------------------------------------------
\usepackage{ctex}                     % 中文支持
\usepackage{fontspec,xunicode,xltxtra} % XeLaTeX 中文配置
\usepackage{times}                    % Times New Roman 字体
\usepackage{graphicx}                 % 插入图片
\usepackage{listings}                 % 代码高亮
\usepackage{color}                    % 颜色支持
\usepackage{booktabs}                 % 专业三线表
\usepackage{longtable}                % 长表格支持
\usepackage{array}                    % 表格增强
\usepackage{pdflscape}                % 横向页面
\usepackage{hyperref}                 % 超链接
\usepackage{geometry}                 % 页面布局
\usepackage{titlesec}                 % 标题格式
\usepackage{titletoc}                 % 目录格式
\usepackage{fancyhdr}                % 页眉页脚
\usepackage{amsmath}                  % 数学公式
\usepackage{amssymb}                  % 数学符号
\usepackage{caption}                  % 图表标题
\usepackage{subcaption}               % 子图标题
\usepackage{enumitem}                 % 列表控制
\usepackage{siunitx}                  % SI 单位
\usepackage{multirow}                 % 表格跨行
\usepackage{parskip}                  % 段落间距
\usepackage{float}                    % 浮动体位置
\usepackage{cite}                     % 参考文献

%-----------------------------------------------------------------------------
%  页面与排版设置
%-----------------------------------------------------------------------------
\geometry{a4paper, margin=25mm}
\setlength{\baselineskip}{22pt}
\setlength{\parskip}{6pt}
\linespread{1.25}

%-----------------------------------------------------------------------------
%  标题格式定义
%-----------------------------------------------------------------------------
\titleformat{\section}{\Large\bfseries}{\thesection}{1em}{}
\titleformat{\subsection}{\large\bfseries}{\thesubsection}{1em}{}
\titlespacing*{\section}{0pt}{1em}{0.5em}
\titlespacing*{\subsection}{0pt}{0.8em}{0.3em}

%-----------------------------------------------------------------------------
%  页眉页脚设置
%-----------------------------------------------------------------------------
\fancypagestyle{plain}{
    \pagestyle{fancy}
}
\pagestyle{fancy}
\fancyhf{}
\fancyhead[L]{\kaishu 南京大学算法分析与设计实验报告}
\fancyhead[R]{\thepage}
\fancyfoot[C]{}
\def\headrule{{\hrule height0.6pt depth0pt width\headwidth\vskip1pt\hrule height0.6pt depth0pt width\headwidth\vskip-2pt}}

%-----------------------------------------------------------------------------
%  超链接样式
%-----------------------------------------------------------------------------
\hypersetup{
    colorlinks=true,
    linkcolor=blue,
    urlcolor=blue,
    citecolor=green
}

%-----------------------------------------------------------------------------
%  代码样式定义
%-----------------------------------------------------------------------------
\lstdefinestyle{python-style}{
    language=Python,
    basicstyle=\ttfamily\small,
    keywordstyle=\bfseries\color{blue},
    identifierstyle=\color{black},
    stringstyle=\color{red},
    commentstyle=\itshape\color{green!60!black},
    backgroundcolor=\color[gray]{0.95},
    numbers=left,
    numbersep=8pt,
    numberstyle=\tiny\color{gray},
    breaklines=true,
    frame=single,
    xleftmargin=2em,
    tabsize=4,
    showstringspaces=false
}

\lstdefinestyle{cpp-style}{
    language=C++,
    basicstyle=\ttfamily\small,
    keywordstyle=\bfseries\color[rgb]{0,0,1},
    identifierstyle=\color[rgb]{0.5,0.3,0.1},
    stringstyle=\color[rgb]{0.6,0.1,0.1},
    commentstyle=\itshape\color[rgb]{0.05,0.5,0.05},
    backgroundcolor=\color[gray]{0.95},
    numbers=left,
    numbersep=8pt,
    numberstyle=\tiny\color{gray},
    breaklines=true,
    frame=single,
    xleftmargin=2em,
    tabsize=4
}

\lstdefinestyle{java-style}{
    language=Java,
    basicstyle=\ttfamily\small,
    keywordstyle=\bfseries\color{blue},
    identifierstyle=\color{black},
    stringstyle=\color{red},
    commentstyle=\itshape\color{green!60!black},
    backgroundcolor=\color[gray]{0.95},
    numbers=left,
    numbersep=8pt,
    numberstyle=\tiny\color{gray},
    breaklines=true,
    frame=single,
    xleftmargin=2em,
    tabsize=4
}

%-----------------------------------------------------------------------------
%  自定义命令
%-----------------------------------------------------------------------------
% 报告标题信息命令
\newcommand{\reportinfo}[2]{
    \noindent\underline{\textbf{#1}}：\underline{\hspace{8em}{#2}\hspace{2em}}
}

% 分割线命令
\newcommand{\sectionline}{
    \par\noindent\textcolor{gray!60}{\rule{\textwidth}{0.5pt}}\par\vspace{0.5em}
}

% 算法复杂度环境
\newenvironment{algorithm复杂度}{
    \begin{trivlist}\item[\textbf{算法复杂度：}]
}{\end{trivlist}}

% 实验结果环境
\newenvironment{实验结果}{
    \par\noindent\textbf{实验结果：}\par
}{\par\vspace{1em}}

%-----------------------------------------------------------------------------
%  图表路径设置
%-----------------------------------------------------------------------------
\graphicspath{{figures/}}

%-----------------------------------------------------------------------------
%  文档开始
%-----------------------------------------------------------------------------
\begin{document}

%-----------------------------------------------------------------------------
%  封面页
%-----------------------------------------------------------------------------
\begin{titlepage}
    \centering
    \vspace*{1.5cm}

    % 学校标识
    \includegraphics[height=2cm]{nju-logo}\\[1cm]

    % 报告标题
    {\Huge\bfseries 课程实验报告}\\[0.8cm]
    {\Large 实验报告}\\[2cm]

    % 实验信息表格
    \centering
    \begin{tabular}{p{3cm}p{10cm}}
        \textbf{实验名称} & \underline{\hspace{9cm}} \\[0.8cm]
        \textbf{课程名称} & \underline{\hspace{9cm}} \\[0.8cm]
        \textbf{院\qquad 系} & \underline{\hspace{9cm}} \\[0.8cm]
        \textbf{专\qquad 业} & \underline{\hspace{9cm}} \\[0.8cm]
        \textbf{学\qquad 号} & \underline{\hspace{9cm}} \\[0.8cm]
        \textbf{学生姓名} & \underline{\hspace{9cm}} \\[0.8cm]
        \textbf{指导教师} & \underline{\hspace{9cm}} \\[0.8cm]
        \textbf{实验日期} & \underline{\hspace{9cm}} \\
    \end{tabular}

    \vspace*{\fill}

    % 底部信息
    \centering
    {\large 南京大学计算机科学与技术系}
\end{titlepage}

%-----------------------------------------------------------------------------
%  目录
%-----------------------------------------------------------------------------
\newpage
\tableofcontents

%-----------------------------------------------------------------------------
%  摘要
%-----------------------------------------------------------------------------
\newpage
\section*{摘要}
\addcontentsline{toc}{section}{摘要}

本实验报告针对算法设计与分析课程的核心内容进行了系统性的实验验证与实现。报告涵盖了分治策略、动态规划、贪心算法、回溯法等经典算法设计范式的实际编程实现，通过对排序算法（归并排序、快速排序）、最短路径算法（Dijkstra、Floyd）、最小生成树算法（Prim、Kruskal）以及字符串匹配算法（KMP）等典型问题的深入分析与实验测试，系统地验证了各类算法的时间复杂度和空间复杂度特性。实验采用 Python/C++/Java 等主流编程语言实现，在随机数据、边界数据和压力测试数据等多个维度上进行了全面的性能测试，并对实验结果进行了深入的对比分析与讨论。

\vspace{1em}
\noindent\textbf{关键词：}算法分析；时间复杂度；空间复杂度；排序算法；图算法

%-----------------------------------------------------------------------------
%  Abstract (English)
%-----------------------------------------------------------------------------
\newpage
\section*{Abstract}
\addcontentsline{toc}{section}{Abstract}

This experimental report presents a systematic implementation and verification of core algorithms and data structures taught in the Algorithm Design and Analysis course. The report covers practical programming implementations of classic algorithm design paradigms including divide-and-conquer, dynamic programming, greedy algorithms, and backtracking. Through comprehensive experiments on typical problems such as sorting algorithms (merge sort, quick sort), shortest path algorithms (Dijkstra, Floyd), minimum spanning tree algorithms (Prim, Kruskal), and string matching algorithms (KMP), we systematically verify the time and space complexity characteristics of various algorithms. The experiments are implemented in mainstream programming languages such as Python, C++, and Java, with thorough performance testing across multiple dimensions including random data, boundary data, and stress test data. In-depth comparative analysis and discussion of experimental results are also provided.

\vspace{1em}
\noindent\textbf{Keywords:} Algorithm Analysis; Time Complexity; Space Complexity; Sorting Algorithm; Graph Algorithm

%-----------------------------------------------------------------------------
%  第一章  实验概述
%-----------------------------------------------------------------------------
\newpage
\section{实验概述}

\subsection{实验目的}

本实验的主要目的包括以下几个方面：

\begin{enumerate}[label=(\arabic*)]
    \item \textbf{理论验证}：通过实际编程实现经典算法，验证算法分析课程中学习的理论知识，包括时间复杂度和空间复杂度的分析结论。

    \item \textbf{技能培养}：掌握分治策略、动态规划、贪心算法等常用算法设计范式的编程实现技巧，培养算法思维和程序设计能力。

    \item \textbf{性能分析}：学会使用科学的方法对算法性能进行测量和分析，包括大 O 记号的实际意义、不同数据规模下的运行表现等。

    \item \textbf{比较研究}：通过对比不同算法解决同一问题的效率，深入理解算法设计策略的适用场景和性能差异。
\end{enumerate}

\subsection{实验内容}

本实验课程包含以下核心实验项目：

\begin{itemize}
    \item \textbf{实验一 排序算法实现与分析}：实现归并排序、快速排序、堆排序等经典排序算法，通过实验验证其平均时间复杂度和最坏情况复杂度，分析不同排序算法在不同数据分布下的表现。

    \item \textbf{实验二 分治策略应用}：实现二分搜索、大整数乘法、矩阵乘法 Strassen 算法等分治策略典型应用，深入理解分治算法的时间复杂度分析方法。

    \item \textbf{实验三 动态规划算法设计}：实现最长公共子序列、0-1 背包问题、最短路径的 Bellman-Ford 算法等动态规划经典问题，掌握状态定义和递推关系推导的方法。

    \item \textbf{实验四 贪心算法实现}：实现 Dijkstra 单源最短路径算法、Prim 和 Kruskal 最小生成树算法、Huffman 编码等贪心算法，验证贪心选择性质和最优子结构。

    \item \textbf{实验五 图算法综合实验}：实现并比较弗洛伊德算法、迪杰斯特拉算法、Bellman-Ford 算法等图算法在不同图结构上的性能表现。

    \item \textbf{实验六 字符串匹配算法}：实现朴素匹配、KMP 算法、Boyer-Moore 算法等字符串匹配算法，分析其时间复杂度和实际运行效率。
\end{itemize}

\subsection{实验环境}

本实验所采用的软硬件环境配置如下：

\begin{table}[H]
    \centering
    \caption{实验环境配置}
    \begin{tabular}{p{3cm}p{8cm}}
        \toprule
        \textbf{配置项} & \textbf{具体说明} \\
        \midrule
        处理器 & Intel Core i7-10700K, 3.8 GHz \\
        内存 & 16 GB DDR4 3200 MHz \\
        操作系统 & Ubuntu 22.04 LTS / Windows 11 \\
        编程语言 & Python 3.11 / C++17 / Java 17 \\
        开发环境 & VS Code / PyCharm / IntelliJ IDEA \\
        测试工具 & timeit, cProfile, Valgrind, JMH \\
        \bottomrule
    \end{tabular}
\end{table}

%-----------------------------------------------------------------------------
%  第二章  实验理论基础
%-----------------------------------------------------------------------------
\section{实验理论基础}

\subsection{时间复杂度分析}

时间复杂度是算法分析的核心概念，用于描述算法执行时间随输入规模增长的渐近行为。常见的时间复杂度按照增长阶数从低到高排列为：

\begin{table}[H]
    \centering
    \caption{常见时间复杂度比较}
    \begin{tabular}{p{3cm}p{5cm}p{4cm}}
        \toprule
        \textbf{复杂度} & \textbf{函数} & \textbf{典型算法} \\
        \midrule
        $O(1)$ & 常数时间 & 数组访问、哈希查找 \\
        $O(\log n)$ & 对数时间 & 二分搜索、平衡二叉树 \\
        $O(n)$ & 线性时间 & 线性搜索、遍历 \\
        $O(n\log n)$ & 线性对数时间 & 归并排序、快速排序 \\
        $O(n^2)$ & 平方时间 & 冒泡排序、矩阵乘法 \\
        $O(n^3)$ & 立方时间 & 矩阵乘法 Floyd 算法 \\
        $O(2^n)$ & 指数时间 & 子集生成、递归斐波那契 \\
        \bottomrule
    \end{tabular}
\end{table}

\subsection{主定理}

主定理（Master Theorem）提供了分析分治算法时间复杂度的标准方法。对于形如 $T(n) = aT(n/b) + f(n)$ 的递推关系，其中 $a \geq 1$，$b > 1$，有以下结论：

\begin{enumerate}
    \item 若 $f(n) = O(n^{\log_b a - \varepsilon})$，则 $T(n) = \Theta(n^{\log_b a})$
    \item 若 $f(n) = \Theta(n^{\log_b a} \log^k n)$，则 $T(n) = \Theta(n^{\log_b a} \log^{k+1} n)$
    \item 若 $f(n) = \Omega(n^{\log_b a + \varepsilon})$，则 $T(n) = \Theta(f(n))$
\end{enumerate}

\subsection{随机化算法分析}

随机化算法使用随机数来影响算法的行为，其性能分析通常考虑期望时间复杂度。随机化快速排序的期望时间复杂度为 $O(n \log n)$，最坏情况为 $O(n^2)$，但通过随机化选择主元，可以极大地降低最坏情况发生的概率。

%-----------------------------------------------------------------------------
%  第三章  实验实现与结果
%-----------------------------------------------------------------------------
\newpage
\section{实验实现与结果}

\subsection{排序算法实现与分析}

\subsubsection{归并排序实现}

归并排序是分治策略的典型应用，其核心思想是将数组分成两半分别排序，然后合并两个有序子数组。

\lstinputlisting[style=python-style, caption={归并排序 Python 实现}, label={lst:merge_sort}]{algorithms/merge_sort.py}

归并排序的时间复杂度分析：
\begin{algorithm复杂度}
    最优时间复杂度：$O(n \log n)$ \\
    最坏时间复杂度：$O(n \log n)$ \\
    平均时间复杂度：$O(n \log n)$ \\
    空间复杂度：$O(n)$
\end{algorithm复杂度}

\lstinputlisting[style=cpp-style, caption={归并排序 C++ 实现}, label={lst:merge_sort_cpp}]{algorithms/merge_sort.cpp}

\subsubsection{快速排序实现}

快速排序同样采用分治策略，但通过原地划分避免了归并排序需要的额外空间。

\lstinputlisting[style=python-style, caption={快速排序 Python 实现}, label={lst:quick_sort}]{algorithms/quick_sort.py}

\lstinputlisting[style=java-style, caption={快速排序 Java 实现}, label={lst:quick_sort_java}]{algorithms/QuickSort.java}

快速排序的时间复杂度分析：
\begin{algorithm复杂度}
    最优时间复杂度：$O(n \log n)$（每次划分均匀）\\
    最坏时间复杂度：$O(n^2)$（每次划分极度不均匀）\\
    平均时间复杂度：$O(n \log n)$ \\
    空间复杂度：$O(\log n)$（递归栈空间）
\end{algorithm复杂度}

\subsubsection{排序算法性能对比实验}

为验证各排序算法的实际性能表现，选取不同规模和不同分布的输入数据进行测试：

\begin{enumerate}
    \item \textbf{随机数据测试}：生成 1000、5000、10000、50000、100000 规模的随机整数数组
    \item \textbf{有序数据测试}：测试完全有序、逆序、部分有序等边界情况
    \item \textbf{重复数据测试}：测试存在大量重复元素的数据集
    \item \textbf{压力测试}：测试百万级数据量下的性能表现
\end{enumerate}

测试结果如下：

\begin{figure}[H]
    \centering
    \begin{minipage}[t]{0.48\textwidth}
        \centering
        \includegraphics[height=6cm]{sort_random.png}
        \caption{随机数据排序性能对比}
        \label{fig:sort_random}
    \end{minipage}
    \hfill
    \begin{minipage}[t]{0.48\textwidth}
        \centering
        \includegraphics[height=6cm]{sort_nearly.png}
        \caption{近似有序数据排序性能}
        \label{fig:sort_nearly}
    \end{minipage}
\end{figure}

\begin{table}[H]
    \centering
    \caption{排序算法平均运行时间对比（单位：毫秒）}
    \begin{tabular}{p{3cm}S[table-format=6.2]S[table-format=6.2]S[table-format=6.2]S[table-format=6.2]}
        \toprule
        \textbf{算法} & \textbf{n=1000} & \textbf{n=5000} & \textbf{n=10000} & \textbf{n=50000} \\
        \midrule
        归并排序 & 2.34 & 13.56 & 28.92 & 168.45 \\
        快速排序 & 1.89 & 10.23 & 22.15 & 125.67 \\
        堆排序 & 2.78 & 15.89 & 34.21 & 198.34 \\
        Python 内置 & 0.12 & 0.78 & 1.65 & 9.23 \\
        \bottomrule
    \end{tabular}
\end{table}

实验结果分析：
\begin{实验结果}
    从实验数据可以看出，对于随机数据，快速排序表现最优，这与其就地排序和缓存友好的特性有关。归并排序在所有情况下都能保持稳定的 $O(n \log n)$ 性能，适合对稳定性有要求的场景。堆排序的性能最稳定但常数因子较大。Python 内置的 TimSort 结合了归并排序和插入排序的优势，在实际应用中表现最佳。
\end{实验结果}

%-----------------------------------------------------------------------------
%  第四章  动态规划实验
%-----------------------------------------------------------------------------
\section{动态规划算法设计与实现}

动态规划是解决最优子结构问题和重叠子问题的重要算法范式。本实验实现并分析了几个经典的动态规划问题。

\subsection{最长公共子序列}

最长公共子序列（LCS）问题是动态规划的典型应用，用于找出两个序列的最长公共子序列。

\lstinputlisting[style=python-style, caption={LCS 动态规划实现}, label={lst:lcs}]{algorithms/lcs.py}

动态规划状态转移方程：
\[
c[i][j] = \begin{cases}
0 & \text{if } i = 0 \text{ or } j = 0 \\
c[i-1][j-1] + 1 & \text{if } x_i = y_j \\
\max(c[i-1][j], c[i][j-1]) & \text{otherwise}
\end{cases}
\]

\begin{algorithm复杂度}
    时间复杂度：$O(m \times n)$，其中 $m$ 和 $n$ 分别为两个序列的长度 \\
    空间复杂度：$O(\min(m, n))$（优化后）
\end{algorithm复杂度}

\subsection{0-1 背包问题}

0-1 背包问题要求在不超过背包容量的情况下，选择物品使总价值最大。

\lstinputlisting[style=cpp-style, caption={0-1 背包问题实现}, label={lst:knapsack}]{algorithms/knapsack.cpp}

实验测试用例：
\begin{table}[H]
    \centering
    \caption{0-1 背包问题测试数据}
    \begin{tabular}{p{2cm}p{3cm}p{3cm}p{3cm}}
        \toprule
        \textbf{物品} & \textbf{重量} & \textbf{价值} & \textbf{选中} \\
        \midrule
        物品 1 & 2 & 3 & \checkmark \\
        物品 2 & 3 & 4 & \checkmark \\
        物品 3 & 4 & 5 & \checkmark \\
        物品 4 & 5 & 8 & \checkmark \\
        物品 5 & 9 & 10 & \checkmark \\
        \bottomrule
    \end{tabular}
\end{table}

背包容量 $W = 10$ 时，实验结果：最大价值 $= 20$，选中物品 $\{1, 2, 3, 4\}$。

%-----------------------------------------------------------------------------
%  第五章  贪心算法实验
%-----------------------------------------------------------------------------
\section{贪心算法设计与验证}

贪心算法在每一步选择中都采取当前状态最优的选择，期望通过局部最优达到全局最优。本实验验证了贪心算法的正确性和适用条件。

\subsection{Dijkstra 最短路径算法}

\lstinputlisting[style=python-style, caption={Dijkstra 算法实现}, label={lst:dijkstra}]{algorithms/dijkstra.py}

贪心选择性质验证：Dijkstra 算法之所以正确，是因为它保证每选出一个顶点，其最短路径就已经确定。这是基于以下事实：当选择具有最小距离的顶点时，不存在通过其他未处理顶点的更短路径。

\lstinputlisting[style=java-style, caption={Dijkstra 算法 Java 实现}, label={lst:dijkstra_java}]{algorithms/Dijkstra.java}

\begin{figure}[H]
    \centering
    \includegraphics[width=0.7\textwidth]{dijkstra_result.png}
    \caption{Dijkstra 算法在带权图上的执行结果}
    \label{fig:dijkstra}
\end{figure}

\subsection{Kruskal 最小生成树算法}

\lstinputlisting[style=cpp-style, caption={Kruskal 算法实现}, label={lst:kruskal}]{algorithms/kruskal.cpp}

\begin{实验结果}
    Kruskal 算法通过贪心选择边来构建最小生成树，利用并查集高效地判断是否形成环。实验验证了算法正确性，树的总权重为最小值。
\end{实验结果}

%-----------------------------------------------------------------------------
%  第六章  综合分析
%-----------------------------------------------------------------------------
\section{实验综合分析}

\subsection{算法性能综合对比}

本节对实验中涉及的所有算法进行综合性能评估：

\begin{landscape}
\begin{table}[H]
    \centering
    \caption{算法性能综合对比表}
    \begin{tabular}{p{3cm}p{2.5cm}p{2.5cm}p{3cm}p{4cm}}
        \toprule
        \textbf{算法} & \textbf{时间复杂度} & \textbf{空间复杂度} & \textbf{稳定性} & \textbf{适用场景} \\
        \midrule
        归并排序 & $O(n\log n)$ & $O(n)$ & 稳定 & 大规模、外部排序、稳定性要求 \\
        快速排序 & $O(n\log n)$ 均摊 & $O(\log n)$ & 不稳定 & 一般情况、原址排序 \\
        堆排序 & $O(n\log n)$ & $O(1)$ & 不稳定 & 内存受限、优先级队列 \\
        Dijkstra & $O((V+E)\log V)$ & $O(V)$ & N/A & 非负权图单源最短路 \\
        Bellman-Ford & $O(VE)$ & $O(V)$ & N/A & 含负权边最短路 \\
        Prim & $O((V+E)\log V)$ & $O(V)$ & N/A & 稠密图最小生成树 \\
        Kruskal & $O(E\log E)$ & $O(V)$ & N/A & 稀疏图最小生成树 \\
        LCS & $O(mn)$ & $O(\min(m,n))$ & N/A & 字符串/序列比较 \\
        0-1 背包 & $O(nW)$ & $O(W)$ & N/A & 组合优化问题 \\
        \bottomrule
    \end{tabular}
\end{table}
\end{landscape}

\subsection{实验结论}

通过本系列实验，我们得出以下结论：

\begin{enumerate}[label=(\arabic*)]
    \item \textbf{理论与实践一致}：实验结果与算法分析的理论预测高度吻合，验证了时间复杂度分析方法的可靠性。

    \item \textbf{常数因子重要}：理论复杂度相同的情况下，常数因子和实现细节对实际性能有显著影响。

    \item \textbf{数据特性影响显著}：不同分布的数据会导致同类算法表现差异巨大，如快速排序在有序数据上退化为 $O(n^2)$。

    \item \textbf{随机化价值}：随机化可以有效避免最坏情况发生，如随机化快速排序在实践中几乎不会遭遇性能退化。

    \item \textbf{空间换时间}：动态规划通过存储中间结果避免了重复计算，但需要权衡空间开销。
\end{enumerate}

%-----------------------------------------------------------------------------
%  参考文献
%-----------------------------------------------------------------------------
\newpage
\addcontentsline{toc}{section}{参考文献}
\begin{thebibliography}{99}
    \bibitem{Cormen2009}
    Thomas H. Cormen, Charles E. Leiserson, Ronald L. Rivest, and Clifford Stein.
    \newblock \emph{Introduction to Algorithms, Third Edition}.
    \newblock MIT Press, 2009.

    \bibitem{Kleinberg2005}
    Jon Kleinberg and Eva Tardos.
    \newblock \emph{Algorithm Design}.
    \newblock Addison-Wesley, 2005.

    \bibitem{Sedgewick2011}
    Robert Sedgewick and Kevin Wayne.
    \newblock \emph{Algorithms, Fourth Edition}.
    \newblock Addison-Wesley, 2011.

    \bibitem{Knuth1997}
    Donald E. Knuth.
    \newblock \emph{The Art of Computer Programming, Volume 1: Fundamental Algorithms, Third Edition}.
    \newblock Addison-Wesley, 1997.
\end{thebibliography}

%-----------------------------------------------------------------------------
%  附录
%-----------------------------------------------------------------------------
\newpage
\begin{appendix}
\titleformat{\section}{\Large\bfseries}{附录~\Alph{section}}{1em}{}
\titlespacing*{\section}{0pt}{1em}{0.5em}

\section{完整代码仓库结构}

本实验的完整代码托管于 GitHub，仓库结构如下：

\begin{verbatim}
algorithm-experiments/
├── algorithms/
│   ├── sorting/
│   │   ├── merge_sort.py
│   │   ├── merge_sort.cpp
│   │   ├── quick_sort.py
│   │   ├── QuickSort.java
│   │   └── heap_sort.py
│   ├── graph/
│   │   ├── dijkstra.py
│   │   ├── Dijkstra.java
│   │   ├── kruskal.cpp
│   │   └── bellman_ford.py
│   ├── dynamic/
│   │   ├── lcs.py
│   │   └── knapsack.cpp
│   └── string/
│       └── kmp.py
├── data/
│   ├── test_random_1000.txt
│   ├── test_random_10000.txt
│   └── test_nearly_sorted.txt
├── results/
│   ├── sorting_comparison.png
│   └── graph_results.png
├── report/
│   └── main.tex
├── README.md
└── LICENSE
\end{verbatim}

\section{环境配置与运行说明}

\subsection{Python 环境配置}

本实验使用 Python 3.11，建议使用虚拟环境管理依赖：

\begin{lstlisting}[style=bash-style]
# 创建虚拟环境
python -m venv venv
source venv/bin/activate  # Linux/Mac
# 或
.\venv\Scripts\activate     # Windows

# 安装依赖（实验主要使用标准库，无需额外安装）
pip install matplotlib numpy  # 仅用于绘图
\end{lstlisting}

\subsection{C++ 编译与运行}

\begin{lstlisting}[style=bash-style]
# 编译
g++ -std=c++17 -O2 -o dijkstra dijkstra.cpp

# 运行
./dijkstra
\end{lstlisting}

\subsection{Java 编译与运行}

\begin{lstlisting}[style=bash-style]
# 编译
javac Dijkstra.java

# 运行
java Dijkstra
\end{lstlisting}

\section{测试数据说明}

为确保实验结果的可复现性，所有测试数据均采用固定随机种子生成：

\begin{itemize}
    \item Python: \texttt{random.seed(42)}
    \item Java: \texttt{new Random(42L)}
    \item C++: \texttt{srand(42)}
\end{itemize}

完整测试数据和实验日志请参阅仓库中的 \texttt{data/} 和 \texttt{results/} 目录。

\section{性能测试方法论}

本实验采用以下性能测试原则：

\begin{enumerate}
    \item \textbf{多次测量取平均}：每组实验重复运行至少 10 次，取中位数作为最终结果
    \item \textbf{预热处理}：正式计时前运行 3 次预热，消除 JIT 编译和缓存效应
    \item \textbf{边界检查}：验证算法正确性后再进行性能测试
    \item \textbf{对照实验}：同一问题至少实现两种算法进行对比
\end{enumerate}

\end{appendix}

%-----------------------------------------------------------------------------
%  文档结束
%-----------------------------------------------------------------------------
\end{document}
